\documentclass[11pt]{article}
\usepackage{amsmath,amsthm,amssymb,mathtools}
\usepackage[margin=1in]{geometry}
\usepackage{hyperref}
\usepackage{booktabs}
\usepackage{enumitem}

\newtheorem{theorem}{Theorem}
\newtheorem{lemma}[theorem]{Lemma}
\newtheorem{proposition}[theorem]{Proposition}
\newtheorem{corollary}[theorem]{Corollary}
\newtheorem{conjecture}[theorem]{Conjecture}
\theoremstyle{definition}
\newtheorem{definition}[theorem]{Definition}
\newtheorem{remark}[theorem]{Remark}
\newtheorem{example}[theorem]{Example}

\DeclareMathOperator{\im}{im}
\DeclareMathOperator{\rank}{rank}
\DeclareMathOperator{\supp}{supp}
\DeclareMathOperator{\dep}{dep}
\DeclareMathOperator{\ord}{ord}
\newcommand{\Vlam}{V^\lambda}
\newcommand{\Ep}[1]{E^{+}_{#1}}
\newcommand{\Em}[1]{E^{-}_{#1}}
\newcommand{\Pq}{P_{q}}
\newcommand{\Pm}{P_{-1}}
\newcommand{\Pdel}{\Pi_{\mathrm{del}}}
\newcommand{\St}{S_{\mathrm T}}
\newcommand{\Sr}{S_{\mathrm R}}
\newcommand{\ddel}{d_{\mathrm{del}}}

\title{A deletion bridge for the off--hook order law:\\
split--independence and a rank--one collapse reduce $\ord=D$\\
to a single--vector reach bound\\
\large (the family $\lambda=(2,2,1^m)$, uniformly in $m$)}
\author{Clio}
\date{31 May 2026 (prove session)}

\begin{document}
\maketitle

\begin{abstract}
The off--hook order law asserts $\ord_{x=q^2}Z_\lambda=D:=\binom{m}{2}$ for $\lambda=(2,2,1^m)$,
equivalently the \emph{split--independent per--subset vanishing} $M(S)\equiv0$ on $\Vlam$ for
every insertion subset $S$ with $|S|<D$. Prior work proved the lower bound $\ord\ge\tau$
($\tau=m-1$) but stalled on the split: the per--block deficit count $g(t)=t+\binom{\tau-t+1}{2}$
is loose ($<D$) whenever the number $t=|\St|$ of tail insertions is positive, and computation
showed the missing killing comes from the \emph{placement} of insertions, not their count. We
remove the split entirely. Writing each inserted projector as $\Pm=I-\Pq$ and expanding the
ordered product, Möbius inversion on the Boolean lattice of word positions yields the
\textbf{deletion bridge}
\[
   M(S)=\sum_{U\subseteq S}(-1)^{|S|-|U|}\,\Pdel(U),\qquad
   \Pdel(U)=\sum_{V\subseteq U}M(V),
\]
where $\Pdel(U)$ is the all--$\Pq$ staircase $\Omega$ with the factors in $U$ \emph{deleted}
(set to $I$). Two consequences are immediate and uniform in $m$: (i) $\ord=\ddel:=\min\{|U|:
\Pdel(U)\neq0\}$; and (ii) \emph{the split--independent per--subset order law is equivalent to
the single split--free statement} $\Pdel(U)=0$ for all $|U|<D$ --- because each $M(S)$ with
$|S|<D$ is a $\pm$--combination of $\Pdel(U)$ with $|U|\le|S|<D$. This dissolves the
intermediate--split and heavy--tail regimes that obstructed the insertion approach. We then show
computationally (exhaustively for $m=3$, by sampling for $m=4$) that for $|U|\le D$ the deletion
operator \emph{collapses through the two longest blocks to rank $\le1$}: the tail sends $\Vlam$
to a line $\langle w\rangle$, and the prefix is a single--vector walk on $w$. The deletion order
law thereby reduces to a \textbf{single--vector reach bound}, the exact off--hook twin of the
unconditional hook chip--firing bound $\binom{m}{2}$ proved for $(2,1^m)$. We isolate the
remaining crux --- a cost--to--go inequality $\Psi(w)\ge D-|\St|$ on the deposited line --- and
show by example that it has no support--only potential (the obstruction is the same
no--cancellation phenomenon as the hook $G_{\binom m2}\neq0$ gap). All identities are verified
by exact rational arithmetic ($q=5/7$, cross--checked $q=2/3$) for $m\le4$.
\end{abstract}

\section{Setup}\label{sec:setup}

We retain the conventions of the 29--30~May deficit papers. $H_n(q)$ is the Iwahori--Hecke
algebra, $T_i^2=(q-1)T_i+q$; on the Hoefsmit seminormal module $\Vlam$ set $S_i:=T_i+1$ and
\[
   \Ep i:=\ker(T_i-q)|_{\Vlam}=\im\Pq^{(i)},\qquad
   \Em i:=\ker(T_i+1)|_{\Vlam}=\im\Pm^{(i)},
\]
with $\Pq^{(i)}=(T_i+1)/(q+1)$ and $\Pm^{(i)}=(qI-T_i)/(q+1)$ the orthogonal eigenprojectors,
$\Pq^{(i)}+\Pm^{(i)}=I$.

\paragraph{The staircase word and the per--subset operators.}
For $\lambda=(2,2,1^m)$ we have $n=m+4$, $N=\binom n2$, $\tau:=\tau(\lambda)=m-1$, and
\[
   D:=\tfrac{\tau(\tau+1)}2=\binom m2 .
\]
Write $B_b:=S_bS_{b-1}\cdots S_1$ for the descending block of length $b$ (generators
$1,\dots,b$), $b=1,\dots,n-1$. The staircase word is the concatenation $B_1B_2\cdots B_{n-1}$;
let $(i_1,i_2,\dots,i_N)$ be its letter sequence. We multiply projector choices left to right
as matrices, so the rightmost letter acts \emph{first} on a vector and the longest block
$B_{n-1}$ acts first. For a subset $S\subseteq\{1,\dots,N\}$ of word positions define the
\emph{insertion} operator
\[
   M(S):=\prod_{k=1}^{N}A_k,\qquad A_k=\Pm^{(i_k)}\ (k\in S),\quad A_k=\Pq^{(i_k)}\ (k\notin S),
\]
and the \emph{deletion} operator
\[
   \Pdel(U):=\prod_{k=1}^{N}C_k,\qquad C_k=I\ (k\in U),\quad C_k=\Pq^{(i_k)}\ (k\notin U).
\]
Thus $\Pdel(\varnothing)=\prod_k\Pq^{(i_k)}=:\Omega$ is the all--$\Pq$ degenerate staircase, and
$\Pdel(U)$ is $\Omega$ with the factors indexed by $U$ removed. As established (operator
reformulation of the order law),
\[
   \ord_{x=q^2}Z_\lambda=\min\{|S|:M(S)\neq0 \text{ on } \Vlam\}=:\ord .
\]
The order law is the assertion $\ord=D$, and its operator form is the
\emph{split--independent per--subset vanishing}: $M(S)\equiv0$ for all $|S|<D$, with no
condition on how $S$ distributes along the word.

\paragraph{Depth.} For $W\subseteq\Vlam$ put
$\dep(W):=\max\{d\ge0:W\subseteq\bigcap_{j=1}^d\Em j\}$.

\paragraph{The obstruction we remove.} The prior insertion analysis splits $S=\Sr\sqcup\St$ into
prefix ($B_1\cdots B_{n-3}$) and tail ($B_{n-2}B_{n-1}$) positions and proves, for $|\St|<\tau$,
the deficit bound $M(S)\neq0\Rightarrow|S|\ge g(|\St|)$ with
$g(t)=t+\binom{\tau-t+1}2$. Since $g(0)=D$ but $g(t)<D$ for $t\ge1$, this gives $\ord\ge\tau$ but
not $\ord=D$; the residual is exactly the placement--sensitive ``intermediate split''
$1\le|\St|<\tau$ together with the heavy--tail regime $|\St|\ge\tau$.

\section{The deletion bridge}\label{sec:bridge}

\begin{theorem}[Deletion bridge]\label{thm:bridge}
As operators on $\Vlam$ (indeed on any module), for every $S\subseteq\{1,\dots,N\}$,
\begin{equation}\label{eq:mobius}
   M(S)=\sum_{U\subseteq S}(-1)^{|S|-|U|}\,\Pdel(U),\qquad\text{and inversely}\qquad
   \Pdel(U)=\sum_{V\subseteq U}M(V).
\end{equation}
\end{theorem}

\begin{proof}
Fix $S$. In the ordered product defining $M(S)$, the factor at a position $k\in S$ is
$\Pm^{(i_k)}=I-\Pq^{(i_k)}$, and at $k\notin S$ it is $\Pq^{(i_k)}$. Expanding each factor
$I-\Pq^{(i_k)}$ ($k\in S$) as a sum of two terms, and keeping every position in its original
slot (no reordering), we obtain
\[
   M(S)=\prod_{k\in S}\bigl(I-\Pq^{(i_k)}\bigr)\!\!\prod_{k\notin S}\!\!\Pq^{(i_k)}
       =\sum_{U\subseteq S}\Bigl(\prod_{k\in U}I\Bigr)\Bigl(\prod_{k\in S\setminus U}(-\Pq^{(i_k)})\Bigr)
        \prod_{k\notin S}\Pq^{(i_k)} ,
\]
where for each $U\subseteq S$ we chose the $I$--term at the positions of $U$ and the
$-\Pq^{(i_k)}$--term at the positions of $S\setminus U$. The sign is $(-1)^{|S\setminus U|}=
(-1)^{|S|-|U|}$, and the surviving operator has $\Pq$ at every position except $U$ (where it has
$I$), i.e.\ it is $\Pdel(U)$. This is the first identity. The second is Möbius inversion on the
Boolean lattice: $\sum_{V\subseteq U}M(V)=\sum_{V\subseteq U}\sum_{W\subseteq V}(-1)^{|V|-|W|}
\Pdel(W)=\sum_{W\subseteq U}\Pdel(W)\sum_{V:W\subseteq V\subseteq U}(-1)^{|V|-|W|}
=\Pdel(U)$, the inner alternating sum vanishing unless $W=U$.
\end{proof}

\begin{corollary}[Order law $=$ deletion order law]\label{cor:ord}
With $\ddel:=\min\{|U|:\Pdel(U)\neq0\}$, one has $\ord=\ddel$.
\end{corollary}

\begin{proof}
If $M(S)\neq0$ then by \eqref{eq:mobius} some $\Pdel(U)\neq0$ with $U\subseteq S$, so
$\ddel\le|U|\le|S|$; minimising over $S$ gives $\ddel\le\ord$. If $\Pdel(U)\neq0$ then some
$M(V)\neq0$ with $V\subseteq U$, so $\ord\le|V|\le|U|$; minimising gives $\ord\le\ddel$.
\end{proof}

\begin{corollary}[Split--independence for free]\label{cor:split}
The following are equivalent, uniformly in $m$:
\begin{enumerate}[topsep=0.2em,itemsep=0.1em,label=\textup{(\alph*)}]
\item $M(S)\equiv0$ on $\Vlam$ for every insertion subset $S$ with $|S|<D$ \textup{(the
split--independent per--subset order law)};
\item $\Pdel(U)\equiv0$ on $\Vlam$ for every deletion subset $U$ with $|U|<D$ \textup{(the
split--free deletion order law)}.
\end{enumerate}
\end{corollary}

\begin{proof}
(b)$\Rightarrow$(a): if $|S|<D$ then every $U\subseteq S$ has $|U|\le|S|<D$, so each term in
$M(S)=\sum_{U\subseteq S}(-1)^{|S|-|U|}\Pdel(U)$ vanishes, whence $M(S)=0$. (a)$\Rightarrow$(b):
symmetric, via $\Pdel(U)=\sum_{V\subseteq U}M(V)$.
\end{proof}

\begin{remark}[Why this is the right move]
Statement (a) carries the entire split apparatus --- the cases $\St=\varnothing$, $1\le|\St|<
\tau$, $|\St|\ge\tau$ all had to be handled separately, and the middle case was where the
count bound $g(t)$ went slack. Statement (b) has \emph{no split}: $U$ is an arbitrary subset of
the whole word. Corollary~\ref{cor:split} says the split was never essential --- it was an
artefact of writing the order law with the projectors $\Pm$ rather than with $I-\Pq$. The
honest object is the single quantity $\ddel$, the minimum number of factors one must delete from
the degenerate staircase $\Omega$ to revive it on $\Vlam$.
\end{remark}

\begin{corollary}[Recovering $\ord\ge\tau$, split--free]\label{cor:tau}
$\ddel\ge\tau$, hence $\ord\ge\tau$; equivalently $\Pdel(U)\equiv0$ for $|U|<\tau$.
\end{corollary}

\begin{proof}
By Corollary~\ref{cor:ord} and the proved bound $\ord\ge\tau$. Alternatively, directly: the
deletion analogue of the deficit lemmas applies verbatim with $I$ in place of $\Pm$. A block
$O_b$ acting on a depth--$d$ subspace $W$ kills it unless its first $d$ letters are deleted
(each $\Pq^{(j)}$, $j\le d$, annihilates $W\subseteq\Em j$, while $I$ fixes it), and the
remaining letters --- generators $\ge d+1$ --- commute past generators $\le d-1$ (axial gap
$\ge2$) and preserve depth $d-1$. Iterating over the prefix blocks below the tail's deposited
depth gives the same arithmetic series, with $\ge\tau$ deletions forced before extinction can be
avoided. \end{proof}

\section{Rank--one collapse: the deletion law is a single--vector walk}\label{sec:collapse}

By Corollary~\ref{cor:split} it remains to study $\Pdel(U)$ for $|U|<D$, an arbitrary subset of
the whole word. The decisive structural fact is that for few deletions the operator factors
through a line.

\begin{proposition}[Rank--one collapse, computational]\label{prop:collapse}
For $\lambda=(2,2,1^m)$ and every $U$ with $|U|\le D$, the running image after the two longest
blocks $B_{n-1}B_{n-2}$ (the tail) has rank $\le1$:
\[
   \rank\bigl(O^U_{n-2}O^U_{n-1}\,\Vlam\bigr)\le1,
\]
where $O^U_b$ is block $B_b$ with its $U$--positions deleted. Hence for $|U|\le D$ the prefix
$B_1\cdots B_{n-3}$ acts on a line, and $\Pdel(U)\neq0$ iff the prefix walk keeps that line
alive.
\end{proposition}

\noindent\emph{Status.} Verified by exact arithmetic: exhaustively over all $U$ with $|U|\le D$
for $m=3$ (where the all--$\Pq$ rank profile of $(B_{n-1},\dots,B_1)$ is $(4,1,0,0,0,0)$), and by
random sampling for $m=4$ (profile $(5,1,0,\dots)$). We expect Proposition~\ref{prop:collapse}
$m$--uniformly from the rank--one pinch (for $|a-b|\ge2$, $\Pq^{(a)}\Pq^{(b)}$ is the projector
onto a $1$--dimensional joint $q$--eigenspace, proved for the whole family) together with the
forward backbone collapse (the all--$\Pq$ prefix has image a single seminormal line, proved
$m$--uniformly). A clean $m$--uniform proof that few deletions cannot defeat enough of these
pinches is the natural next lemma; we treat the collapse as a hypothesis for the reduction below.

\begin{corollary}[Reduction to a single--vector reach bound]\label{cor:reduce}
Assume Proposition~\ref{prop:collapse}. Split $U=U_{\mathrm T}\sqcup U_{\mathrm R}$ into tail and
prefix positions, and for a deposited line $\langle w\rangle=O^{U_{\mathrm T}}_{n-2}
O^{U_{\mathrm T}}_{n-1}\Vlam$ let
\[
   \Psi(w):=\min\{|U_{\mathrm R}|:\text{the prefix with deletions }U_{\mathrm R}\text{ keeps }
   \langle w\rangle\text{ nonzero}\}
\]
be the cost--to--go. Then the deletion order law $\ddel\ge D$ is equivalent to the
\emph{single--vector reach bound}
\begin{equation}\label{eq:reach}
   \Psi(w)\ \ge\ D-|U_{\mathrm T}|\qquad\text{for every tail--deposited line }w .
\end{equation}
\end{corollary}

\begin{proof}
If $\Pdel(U)\neq0$ with $U=U_{\mathrm T}\sqcup U_{\mathrm R}$, the prefix keeps the deposited line
alive, so $|U_{\mathrm R}|\ge\Psi(w)\ge D-|U_{\mathrm T}|$ by \eqref{eq:reach}, giving
$|U|\ge D$. Conversely, taking $U_{\mathrm T}=\varnothing$, \eqref{eq:reach} reads
$\Psi(w_0)\ge D$ for the all--$\Pq$ deposited line $w_0$, which is exactly $\ddel\ge D$ on the
all--$\Pq$--tail slice; the displayed inequality for $|U_{\mathrm T}|\ge1$ supplies the
remaining slices.
\end{proof}

\noindent The inequality \eqref{eq:reach} is the off--hook twin of the unconditional chip--firing
lower bound for the hook $(2,1^m)$, where $\Pq^{(i)}$ is rank one on $V^{(2,1^m)}$, the trace
factorises into scalar sandwiches, and a reach principle (a descending run raises the reachable
maximum by at most one) gives $\ord\ge\binom m2$ with no positivity input. After the collapse of
Proposition~\ref{prop:collapse} the off--hook prefix is likewise a single--vector walk, and we
expect the same reach principle.

\section{The crux, and why no support--only potential closes it}\label{sec:crux}

The tautological route to \eqref{eq:reach} fails: the cost--to--go $\Psi$ satisfies the
super--triangle inequality $\Psi(O_b^d\langle v\rangle)\ge\Psi(\langle v\rangle)-|d|$ for any
block deletion pattern $d$ (immediate from its definition as a minimum), but this gives no
\emph{lower} bound on $\Psi$ of the deposited line without an independent handle. One therefore
seeks a \emph{closed--form potential} $\varphi$ on lines with
\[
   \text{(i) } \varphi(w)\ge D-|U_{\mathrm T}|\ \text{ on tail--deposited lines},\qquad
   \text{(ii) } \varphi(O_b^d\langle v\rangle)\ge\varphi(\langle v\rangle)-|d| ,
\]
so that summing (ii) over the prefix blocks yields $|U_{\mathrm R}|\ge\varphi(w)\ge
D-|U_{\mathrm T}|$.

\paragraph{The natural candidate, and its failure.} Each seminormal basis vector of $\Vlam$
corresponds to its column--$1$ pair $(a,b)$, $a<b$, the entries of the two cells $(0,1),(1,1)$;
the row--superstandard target line is $(a,b)=(2,4)$. The graph distance to the target,
$\mathrm{ch}(a,b):=(a-2)+(b-4)$, gives a valid \emph{lower} bound
$\Psi(w)\ge\min_{(a,b)\in\supp(w)}\mathrm{ch}(a,b)$ on every line tested ($m=3$). But this
minimum is \emph{not} $\ge D-|U_{\mathrm T}|$: at $m=3$, $|U_{\mathrm T}|=2$, the deposited line
with support $\{(2,4),(2,5),(3,4),(3,5)\}$ has $\min\mathrm{ch}=0$ (the target itself lies in the
support), yet $\Psi=1=D-|U_{\mathrm T}|$. The cheap $(2,4)$--component is present in the support
but \emph{cancels} along the cheapest prefix walk, so the support alone does not determine
$\Psi$.

\paragraph{Identification of the obstruction.} Two distinct residuals meet here. The
\emph{lower} bound \eqref{eq:reach} is obstructed by the off--hook ``merged--arc'' (merged
junction) gap: when few deletions break a shared pinch, a support--based reach principle
\emph{undershoots} the true cost --- exactly as $\min_{(a,b)\in\supp}\mathrm{ch}$ gives $0$ on a
line whose true $\Psi$ is $1$. So a coefficient--blind reach principle does not suffice for the
lower bound; one needs either a refined reach that charges merged pinches correctly, or a
coefficient--aware argument. (Separately, the \emph{equality} $\ord=\binom m2$ rather than $\ge$
is the no--cancellation residual $G_{\binom m2}=\sum_{|S|=\binom m2}M(S)\neq0$, already in hand
from the survivor census.) The single--vector reach bound \eqref{eq:reach} is therefore not a
function of the line's $\Em\bullet$ signature or of its support; we record it as the irreducible
core of the lower bound.

\begin{conjecture}[Single--vector reach bound]\label{conj:reach}
For $\lambda=(2,2,1^m)$ the inequality \eqref{eq:reach} holds: every tail--deposited line $w$
needs at least $D-|U_{\mathrm T}|$ prefix deletions to survive. Equivalently
$\ddel\ge D$, equivalently (Corollary~\ref{cor:split}) the split--independent order law $\ord=D$.
\end{conjecture}

\section{Verification}\label{sec:verify}

All checks use the Hoefsmit seminormal machinery; structural identities are exact rational
($q=5/7$, cross--checked $q=2/3$), bulk screens use a generic rational $q$ to exclude accidental
cancellation. Scripts in \texttt{\~{}/projects/scratch/2026-05-31-intermediate-split/}.

\begin{itemize}[topsep=0.2em,itemsep=0.1em]
\item \emph{Bridge \eqref{eq:mobius} as matrices.} Both identities hold exactly for randomly
chosen $S,U$ at $m=3$ and $m=4$.
\item \emph{Deletion order law $\ddel=D$.} $m=2$: $\ddel=1=D$. $m=3$: all $|U|\le2$ deletions
give $\Pdel(U)=0$ exactly ($1+21+210$ subsets), first survivor at $|U|=3=D$. $m=4$: all $|U|\le5$
give $0$ (generic--$q$ screen with early zero short--circuit), survivor at $|U|=6=D$.
\item \emph{Rank--one collapse.} For all $U$ with $|U|\le D$ at $m=3$ the running rank after the
tail is $\le1$; likewise on $60$ random subsets at $m=4$. The all--$\Pq$ rank profiles are
$(4,1,0,\dots)$ ($m=3$) and $(5,1,0,\dots)$ ($m=4$).
\item \emph{Cost--to--go.} For every tail--deposited line at $m=3$, $\Psi(w)\ge D-|U_{\mathrm
T}|$ holds (minimum over deposits at fixed $|U_{\mathrm T}|$ is attained with equality, giving the
$D-|U_{\mathrm T}|$ survivors); and $\Psi$ is not a function of the $\Em\bullet$ signature
(signature $(1,2)$ occurs with $\Psi\in\{3,5\}$).
\end{itemize}

\begin{center}
\renewcommand{\arraystretch}{1.15}
\begin{tabular}{c|c|c|c|c|c}
\toprule
$\lambda$ & $m$ & $\tau$ & $D$ & $\ddel$ (computed) & first survivor $|U|$\\
\midrule
$(2,2,1,1)$     & 2 & 1 & 1 & $1$ & $1$\\
$(2,2,1,1,1)$   & 3 & 2 & 3 & $3$ & $3$\\
$(2,2,1,1,1,1)$ & 4 & 3 & 6 & $6$ & $6$\\
\bottomrule
\end{tabular}
\end{center}

\section{Status}\label{sec:status}

\paragraph{Proved (uniformly in $m$).}
\begin{itemize}[topsep=0.2em,itemsep=0.1em]
\item Theorem~\ref{thm:bridge} (the deletion bridge) --- an elementary but exact operator
identity, valid on any module.
\item Corollary~\ref{cor:ord}: $\ord=\ddel$.
\item Corollary~\ref{cor:split}: the split--independent per--subset order law is \emph{equivalent}
to the split--free deletion order law $\Pdel(U)\equiv0$ for $|U|<D$. This is the structural
payoff: it eliminates the intermediate--split and heavy--tail casework.
\item Corollary~\ref{cor:tau}: $\ddel\ge\tau$ (hence $\ord\ge\tau$), now split--free.
\end{itemize}

\paragraph{Reduced (modulo Proposition~\ref{prop:collapse}).}
The deletion order law $\ddel\ge D$ is equivalent to the single--vector reach bound
\eqref{eq:reach} (Corollary~\ref{cor:reduce}). Proposition~\ref{prop:collapse} (the rank--one
collapse through the tail) is verified for $m\le4$ and expected $m$--uniformly from the
rank--one pinch and the forward backbone collapse.

\paragraph{Open (the crux).}
Conjecture~\ref{conj:reach}: the single--vector reach bound $\Psi(w)\ge D-|U_{\mathrm T}|$. We
have shown it is not a function of the line's support or eigenspace signature; the obstruction to
this \emph{lower} bound is the off--hook merged--arc gap (a support--based reach undershoots when
a shared pinch is broken). The reduction to a \emph{single} split--free inequality on a
\emph{single} vector --- the off--hook twin of the proved unconditional hook chip--firing bound
$\binom m2$ --- is the contribution of this note; closing it is the remaining work.
\end{document}
